\documentclass[a4paper,12pt]{article}
\usepackage[T2A]{fontenc}
\usepackage[utf8]{inputenc}
\usepackage[russian]{babel}
\usepackage[a4paper,margin=22mm,headheight=15pt]{geometry}
\usepackage{amsmath,amssymb,booktabs,tabularx,xcolor,tikz,fancyhdr}
\usepackage[colorlinks=true,linkcolor=teal!70!black]{hyperref}
\definecolor{accent}{HTML}{008080}
\definecolor{darkaccent}{HTML}{004D4D}
\setlength{\parindent}{0pt}
\setlength{\parskip}{5pt}
\setlength{\abovedisplayskip}{5pt}
\setlength{\belowdisplayskip}{5pt}
\renewcommand{\arraystretch}{1.35}
\pagestyle{fancy}
\fancyhf{}
\fancyhead[L]{\small Задачи на движение}
\fancyhead[R]{\small Ярослав Гаврилов}
\fancyfoot[C]{\thepage}
\renewcommand{\headrulewidth}{0.3pt}
\newcommand{\task}[2]{\clearpage\phantomsection\label{task:#1}\section*{\textcolor{darkaccent}{Задание №#1. #2}}}
\newcommand{\parthead}[1]{\par\textbf{#1.}\ }
\newcommand{\error}[1]{\par\textcolor{red!75!black}{\textbf{Ошибка:}} #1\par}
\newcommand{\idea}[1]{\par\textcolor{blue!60!black}{\textbf{Ключевая идея:}} #1\par}
\newcommand{\result}[1]{\par\textcolor{darkaccent}{\textbf{Итог:}} #1\par}
\newcommand{\practice}[1]{\par\medskip\textcolor{accent}{\textbf{Задача на закрепление}}\par #1}
\begin{document}
\begin{titlepage}
\centering
\vspace*{25mm}
{\Large\color{darkaccent} МАТЕМАТИКА\par}
\vspace{14mm}
{\huge Ревью домашней работы\par}
\vspace{10mm}
{\LARGE\color{accent} Задачи на движение\par}
\vspace{12mm}
{\large Средняя скорость, встреча, обгон и течение\par}
\vfill
{\large Ярослав Гаврилов\par}
\vspace{8mm}
26 ноября 2026 г.\par
\vspace{16mm}
{\small Лёвин Артём Александрович\par эксклюзивно для Ярослава Гаврилова\par}
\vspace{14mm}
\begin{tikzpicture}
\draw[accent,line width=.8pt] (0,0) .. controls (2,0.8) and (4,-0.8) .. (6,0) .. controls (8,0.8) and (10,-0.8) .. (12,0);
\end{tikzpicture}
\end{titlepage}
\section*{\textcolor{darkaccent}{Сводная таблица}}
Нажмите на номер задания, чтобы перейти к подробному разбору. В таблице приведены только задания, требующие доработки.
\par\medskip
\noindent\begin{tabularx}{\linewidth}{@{}c >{\raggedright\arraybackslash}p{28mm} >{\raggedright\arraybackslash}X >{\raggedright\arraybackslash}p{29mm} >{\raggedright\arraybackslash}p{30mm}@{}}
\toprule
№ & Тема & Суть ошибки или недочёта & Ваш ответ & Правильный ответ\\
\midrule
\hyperref[task:2]{2} & Средняя скорость & Не выполнено сравнение & 72 км/ч & 72 км/ч, на 3 км/ч меньше\\
\hyperref[task:6]{6} & Обгон пешехода & Нет объяснения разности скоростей & 570 (м) & 570 м с обоснованием\\
\hyperref[task:7]{7} & Встречные поезда & Смешаны единицы длины и скорости; неверная единица ответа & 15 минут & 15 с\\
\hyperref[task:8]{8} & Полный обгон & Преждевременное округление & 50 & 50 с, точно\\
\hyperref[task:9]{9} & Движение по реке & Не обоснован отбор корня & $\pm9$; итоговая цифра неясна & 9 км/ч\\
\bottomrule
\end{tabularx}
\clearpage
\thispagestyle{empty}
\vspace*{\fill}
\begin{center}{\Huge\color{darkaccent} Разбор ошибок}\end{center}
\vspace*{\fill}
\task{2}{Две половины маршрута}
\parthead{Текст задания} Поездка составила 180 км. Первые 90 км автомобиль проехал со скоростью 60 км/ч, остальные 90 км --- со скоростью 90 км/ч. Остановок не было. Найдите среднюю скорость и сравните её с результатом предыдущей задачи.
\parthead{Ваше решение} В таблице найдены времена $90:60=1{,}5$ ч и $90:90=1$ ч. Затем записано $v_{\text{ср}}=180:2{,}5=72$.
\parthead{Ваш ответ} 72 км/ч (единица следует из расчёта).
\error{Не выполнена вторая часть условия. Численный расчёт верен.}
\parthead{Где именно} Последний правильный шаг --- получение 72 км/ч. Следующий необходимый шаг, сравнение с результатом предыдущей задачи, отсутствует. Это не вычислительная ошибка, а неполный ответ.
\parthead{Почему это важно} При равных расстояниях автомобиль дольше движется с меньшей скоростью. Поэтому среднее арифметическое двух скоростей здесь не является средней скоростью всей поездки.
\idea{Средняя скорость --- весь путь, делённый на всё время. При равных временах и при равных расстояниях получаются разные результаты.}
\parthead{Исправление} Сохраните Ваш расчёт и допишите сравнение: $75-72=3$ км/ч.
\parthead{Полное решение} Время на участках:
\[t_1=\frac{90}{60}=1{,}5\ \text{ч},\qquad t_2=\frac{90}{90}=1\ \text{ч}.\]
Общее время равно $2{,}5$ ч, а общий путь --- 180 км. Следовательно,
\[v_{\text{ср}}=\frac{180}{2{,}5}=72\ \text{км/ч}.\]
При равных временах средняя скорость равна $(60+90):2=75$ км/ч. Значит, при равных расстояниях она на 3 км/ч меньше.
\parthead{Проверка} $72\cdot2{,}5=180$ км; результат находится между 60 и 90 км/ч.
\result{72 км/ч, на 3 км/ч меньше результата предыдущей задачи.}
\practice{Автомобиль проехал первые 120 км со скоростью 60 км/ч, следующие 120 км --- со скоростью 80 км/ч. Найдите среднюю скорость и сравните её со средней скоростью при движении с теми же скоростями в течение равных промежутков времени.}
\task{6}{Поезд обгоняет пешехода}
\parthead{Текст задания} Поезд со скоростью 72 км/ч обгоняет пешехода, идущего в том же направлении со скоростью 3,6 км/ч. Полное прохождение поезда мимо пешехода занимает 30 с. Найдите длину состава в метрах. Объясните выбор разности скоростей.
\parthead{Ваше решение} Записано $\frac{30}{3600}=\frac{x}{72-3{,}6}$, затем $x=0{,}57$ км, или 570 м.
\parthead{Ваш ответ} 570 (м).
\error{Численный ответ верен, но требуемое объяснение выбора разности скоростей отсутствует.}
\parthead{Где именно} Последний правильный шаг --- использование относительной скорости $72-3{,}6$. Не пояснено, почему скорости нужно вычитать, а не складывать.
\parthead{Почему это важно} Оба объекта движутся в одну сторону. За каждую секунду пешеход тоже продвигается вперёд. Поезд выигрывает у него не весь свой путь, а разность двух путей за это время. Когда последний вагон поравняется с пешеходом, относительное перемещение поезда будет равно его длине.
\idea{При обгоне в одном направлении относительная скорость равна разности большей и меньшей скоростей.}
\parthead{Исправление} Допишите это объяснение перед составлением уравнения. Само уравнение менять не нужно.
\parthead{Полное решение} Удобно сразу перейти к метрам и секундам. Поскольку 1 км/ч равен $\frac{1000}{3600}=\frac1{3{,}6}$ м/с, получаем
\[v_{\text{поезда}}=72:3{,}6=20\ \text{м/с},\quad v_{\text{пешехода}}=3{,}6:3{,}6=1\ \text{м/с}.\]
За секунду поезд выигрывает $20-1=19$ м. За 30 секунд:
\[L=19\cdot30=570\ \text{м}.\]
\parthead{Проверка} Поезд проходит 600 м, пешеход --- 30 м. Разность составляет 570 м, то есть длину состава.
\result{570 м. При обгоне скорости вычитаются, потому что оба объекта движутся в одном направлении.}
\practice{Поезд движется со скоростью 90 км/ч и полностью обгоняет пешехода, идущего в том же направлении со скоростью 7,2 км/ч, за 20 секунд. Найдите длину поезда в метрах и объясните выбор относительной скорости.}
\task{7}{Два встречных поезда}
\parthead{Текст задания} По параллельным путям навстречу друг другу движутся поезда длиной 240 и 360 м со скоростями 54 и 90 км/ч соответственно. Сколько секунд проходит от встречи локомотивов до полного расхождения последних вагонов?
\parthead{Ваше решение} В расчёте использованы сумма длин и сумма скоростей: $\frac{240+360}{54+90}=\frac{600}{144}$. Далее записано $4{,}16\cdot3600=15$ минут.
\parthead{Ваш ответ} 15 минут.
\error{Длины в метрах разделены на скорость в километрах в час; затем неверно переведено время.}
\parthead{Где именно} Суммы 600 м и 144 км/ч выбраны правильно. Первый неверный переход --- трактовка их отношения как времени в часах без согласования единиц. Кроме того, $4{,}16\cdot3600=14976$, а не 15.
\parthead{Почему это неверно} В формуле $t=S/v$ единицы должны согласовываться: метры и метры в секунду либо километры и километры в час. Умножение часов на 3600 даёт секунды, не минуты.
\idea{При встречном движении относительная скорость равна сумме скоростей, а для полного расхождения требуется пройти сумму длин.}
\parthead{Исправление} Сохраните обе суммы, но переведите 144 км/ч в м/с: $144:3{,}6=40$ м/с.
\parthead{Полное решение} От встречи локомотивов до расхождения последних вагонов поезда должны взаимно сместиться на $240+360=600$ м. За секунду они сближают соответствующие точки на
\[54:3{,}6+90:3{,}6=15+25=40\ \text{м/с}.\]
Значит, $t=600:40=15$ с. В километрах и часах тот же расчёт выглядит так:
\[t=\frac{0{,}6}{144}\cdot3600=15\ \text{с}.\]
\parthead{Проверка} За 15 с поезда пройдут 225 и 375 м, в сумме 600 м.
\result{15 секунд, а не минут.}
\practice{Два поезда длиной 280 и 420 м движутся навстречу друг другу со скоростями 72 и 54 км/ч. Сколько секунд пройдёт от встречи локомотивов до полного расхождения последних вагонов?}
\task{8}{Обгон одного поезда другим}
\parthead{Текст задания} Поезда длиной 180 и 320 м движутся в одном направлении со скоростями 108 и 72 км/ч соответственно. Найдите время полного обгона: от момента, когда локомотив быстрого поезда поравнялся с последним вагоном медленного, до момента, когда последний вагон быстрого обогнал локомотив медленного.
\parthead{Ваше решение} После первоначальной записи длин в метрах Вы перевели их в километры и получили $\frac{0{,}18+0{,}32}{36}$. Далее записано $0{,}014\cdot3600=50$.
\parthead{Ваш ответ} 50 (с по смыслу условия).
\error{Преждевременное округление и неточный знак равенства. Итоговое число верно.}
\parthead{Где именно} Последний правильный шаг --- дробь $\frac{0{,}5}{36}$ ч. Первый неверный переход --- замена её на $0{,}014$ ч как на точное значение.
\parthead{Почему это неверно} Дробь равна $\frac1{72}$ ч, то есть приблизительно $0{,}013888\ldots$ ч. При округлении до тысячных теряется точность: $0{,}014\cdot3600=50{,}4$, а не 50. Если округление необходимо, следует использовать знак $\approx$; здесь оно вообще не требуется.
\idea{При полном обгоне относительное перемещение равно сумме длин; точную дробь лучше сохранять до последнего действия.}
\parthead{Исправление} Продолжите с верной дроби:
\[t=\frac{0{,}5}{36}\cdot3600=50\ \text{с}.\]
\parthead{Полное решение} Быстрый поезд сначала проходит всю длину медленного относительно него, а затем ещё собственную длину. Поэтому $S_{\text{отн}}=180+320=500$ м. Относительная скорость:
\[v_{\text{отн}}=(108-72):3{,}6=10\ \text{м/с}.\]
Время равно $500:10=50$ с.
\parthead{Проверка} За 50 секунд быстрый поезд выигрывает у медленного ровно $10\cdot50=500$ м.
\result{50 с точно. Сохраняйте дробь либо сразу переводите скорости в м/с.}
\practice{Поезда длиной 210 и 330 м движутся в одном направлении со скоростями 99 и 63 км/ч. Сколько секунд длится полный обгон от встречи локомотива быстрого поезда с последним вагоном медленного до расхождения последнего вагона быстрого с локомотивом медленного?}
\task{9}{Катер на реке}
\parthead{Текст задания} Катер прошёл 24 км против течения и вернулся тем же путём. Скорость течения 3 км/ч. Против течения катер затратил на 2 ч больше. Найдите собственную скорость катера, учитывая, что она больше скорости течения.
\parthead{Ваше решение} Составлено верное уравнение $\frac{24}{x-3}-\frac{24}{x+3}=2$. В незачёркнутой части получено $2(x^2-9)=144$, затем $x^2-9=72$ и $x=\pm9$.
\parthead{Ваш ответ} В решении указаны два корня, $\pm9$; последняя цифра после слова «От.» читается неоднозначно.
\error{Не записано ограничение $x>3$ и не объяснён отбор корня.}
\parthead{Где именно} Ограничение следовало указать ещё при введении переменной. Последний надёжно читаемый правильный результат --- $x=\pm9$ как корни квадратного уравнения. Необходимо явно исключить $-9$. Зачёркнутые попытки не оцениваются.
\parthead{Почему это важно} Собственная скорость положительна. Более того, чтобы катер действительно двигался против течения, его скорость $x-3$ должна быть строго положительной. Поэтому не всякий алгебраический корень подходит по смыслу задачи.
\idea{Перед составлением дробного уравнения фиксируйте ограничения, а после решения проверяйте каждый корень.}
\parthead{Исправление} Из $x=\pm9$ и $x>3$ следует $x=9$.
\parthead{Полное решение} Пусть $x>3$ --- собственная скорость в км/ч. Время против течения равно $24/(x-3)$, по течению --- $24/(x+3)$. Разность времён равна 2 ч. Приведём дроби к общему знаменателю:
\[\frac{24(x+3)-24(x-3)}{(x-3)(x+3)}=2,\qquad \frac{144}{x^2-9}=2.\]
При $x>3$ знаменатель ненулевой. Умножаем обе части на него:
\[2(x^2-9)=144,\quad x^2-9=72,\quad x^2=81.\]
Отсюда $x=9$ или $x=-9$. Второй корень не удовлетворяет $x>3$.
\parthead{Проверка} Против течения: $24:6=4$ ч; по течению: $24:12=2$ ч. Разность равна 2 ч.
\result{9 км/ч. В ответ записывается только допустимый корень.}
\practice{Катер прошёл 24 км против течения и вернулся. Скорость течения 2 км/ч, а путь против течения занял на 3 часа больше. Найдите собственную скорость катера, укажите ограничение и обоснуйте отбор корня.}
\end{document}
